\section{Related Work}\label{related-work}

Personality-adaptive conversational AI sits at the intersection of affective computing, digital mental health interventions, personality-aware dialogue systems, and LLM-based conversational architectures. We synthesize this landscape in two arcs: (i) the evolution from emotion recognition to deployed support tools and personality-aware dialogue, and (ii) the rise of LLMs alongside the resulting orchestration and evaluation challenges.

\subsection{From Affective Computing to Personality-Aware Digital Support}\label{from-affective-computing-to-personality-aware-digital-support}

Early affective computing focused on recognizing transient emotional states from text, speech, and facial expressions {[}12,13{]}, evolving from categorical models toward dimensional representations (valence/arousal). A persistent limitation is the emphasis on short-lived states without incorporating relatively stable individual differences that shape how emotions are processed and how support is received {[}6{]}.

In parallel, commercial companions and research interventions (e.g., Replika, Woebot, ElliQ) demonstrate that conversational agents can deliver scalable support and reduce loneliness-related outcomes {[}21--24{]}. However, many systems remain ``one-size-fits-all,'' relying on fixed interaction strategies or shallow personalization based on preferences rather than stable psychological profiles {[}22--24{]}.

Personality-aware dialogue research suggests that adapting language style and interaction strategies to Big Five traits can improve perceived naturalness and engagement {[}6,30{]}. Yet most approaches rely on static initialization (explicit questionnaires or one-time inference) and do not update as evidence accumulates, limiting long-term consistency and attunement {[}6,30{]}.

To move beyond surface-level stylistic matching, computational systems also require a cognitive-motivational rationale for why traits should alter interaction strategies. The Zurich Model of Social Motivation {[}10,26--29{]} provides this grounding by linking Big Five traits to three motivational domains---security, arousal, and affiliation---offering a principled basis for mapping detected traits to regulation behaviors.

\subsection{LLMs in Mental Health: Capabilities and Evaluation Challenges}\label{llms-in-mental-health-capabilities-and-evaluation-challenges}

LLMs have transformed conversational AI by enabling flexible reasoning, context integration, and fluent natural language generation {[}8,15,16{]}. Instruction-tuned variants---trained to follow human-specified directives via reinforcement learning from human feedback---further improved controllability and alignment {[}64{]}. Recent work demonstrates that LLM-based systems can support therapeutic-style dialogue through dialogue state tracking, cascaded empathy mechanisms, and multimodal extensions {[}15--20{]}. Despite these capabilities, LLM systems often adapt to immediate conversational context rather than maintaining persistent user models, making relatively stable personalization (e.g., personality) an open challenge.

Operationalizing multi-turn personalization with LLMs typically requires orchestration to balance flexibility with behavioral control. PROMISE is a model-driven framework for developing complex language-based interactions that uses state machine modeling concepts---states, transitions, and dynamic prompt orchestration---to improve behavioral control and reproducibility of LLM-based interactions {[}7{]}. The PROMISE technical report details how prompts are composed along hierarchically nested states and transitions (including triggers, guards, and actions), and provides health information system scenarios and an accompanying public codebase implementing these abstractions {[}7{]}.

Evaluation remains a central methodological challenge. Human expert ratings are costly and difficult to scale {[}25{]}, motivating ``LLM-as-judge'' approaches; however, AI-evaluating-AI can introduce circularity and bias concerns {[}55,56{]}. These tensions motivate hybrid strategies where structured automated evaluation is benchmarked against human expert judgment.

\textbf{Synthesis of gaps and contributions}: This literature highlights (i) the need for dynamic (not static) personality modeling, (ii) the need for theory-driven mappings from traits to behavior (e.g., via Zurich Model domains), and (iii) the need for evaluation methods that scale while maintaining credibility. Our study addresses these gaps via a modular detection$\rightarrow$regulation$\rightarrow$evaluation pipeline and a structured LLM-based evaluation framework, supplemented by author-conducted qualitative review. The \emph{Evaluation Matrix}---a structured per-criterion rubric with distinct criterion sets for regulated versus baseline agents (regulated responses measured in five dimensions: detection accuracy, regulation effectiveness, emotional tone, coherence, and psychological need fulfillment)---was conceived in prior work {[}68{]} (``a per-criterion scoring matrix designed for LLM-evaluator assessment of personality-adaptive response quality, using a trinary Yes/Not-Sure/No scale''). Independent multi-rater validation remains essential future work.

